%==============================================================================
% DCC062 - Sistemas Operacionais - Trabalho Teorico (2026.1)
% Tema 1: Android e Fuchsia
% Pontos de responsabilidade: 1, 2 e 4
% Autor: Felipe Lazzarini Cunha
%==============================================================================
\documentclass[12pt,a4paper]{article}

%------------------------------------------------------------------ Codificacao
\usepackage[utf8]{inputenc}
\usepackage[T1]{fontenc}
\usepackage[brazilian]{babel}

%------------------------------------------------------------------ Tipografia
\usepackage{lmodern}
\usepackage{microtype}            % ajuste fino de espacamento
\usepackage{setspace}
\onehalfspacing                   % espacamento 1,5 (padrao academico)

%--------------------------------------------------------------------- Layout
\usepackage[a4paper,left=3cm,right=2cm,top=3cm,bottom=2cm]{geometry}
\usepackage{titlesec}
\usepackage{parskip}              % paragrafos com espaco em vez de recuo
\usepackage{enumitem}
\usepackage{graphicx}
\usepackage{xcolor}
\usepackage{booktabs}             % tabelas profissionais
\usepackage{array}

%------------------------------------------------------------------ Hyperlinks
\definecolor{linkazul}{HTML}{1A5276}
\usepackage[hidelinks]{hyperref}
\hypersetup{
    colorlinks=true,
    linkcolor=linkazul,
    citecolor=linkazul,
    urlcolor=linkazul,
    pdftitle={Trabalho Teorico - Android e Fuchsia},
    pdfauthor={Felipe Lazzarini Cunha}
}
\usepackage[capitalise,brazilian]{cleveref}

%--------------------------------------------------------- Estilo dos titulos
\titleformat{\section}
  {\normalfont\Large\bfseries\color{linkazul}}{\thesection}{1em}{}
\titleformat{\subsection}
  {\normalfont\large\bfseries}{\thesubsection}{1em}{}

%------------------------------------------------------ Cabecalho / rodape
\usepackage{fancyhdr}
\pagestyle{fancy}
\fancyhf{}
\setlength{\headheight}{13.6pt}
\renewcommand{\headrulewidth}{0.4pt}
\lhead{\small DCC062 --- Sistemas Operacionais}
\rhead{\small Tema 1: Android e Fuchsia}
\cfoot{\thepage}

%------------------------------------------------------ Ambiente de "nota"
% Caixa para anotacoes/lembretes de pesquisa (remover na versao final)
\usepackage{mdframed}
\newmdenv[
  linecolor=linkazul,linewidth=0.6pt,
  backgroundcolor=linkazul!4,
  roundcorner=4pt,skipabove=8pt,skipbelow=8pt,
  innertopmargin=6pt,innerbottommargin=6pt
]{nota}

\begin{document}
\pagenumbering{gobble}

%==============================================================================
% CAPA
%==============================================================================
\begin{titlepage}
\centering
\vspace*{1cm}

{\scshape\large Universidade Federal de Juiz de Fora\par}
{\scshape Departamento de Ciencia da Computacao\par}
\vspace{0.4cm}
{\scshape DCC062 --- Sistemas Operacionais --- 2026.1\par}

\vfill

{\color{linkazul}\rule{\linewidth}{0.6pt}}
\vspace{0.6cm}

{\Huge\bfseries Android e Fuchsia\par}
\vspace{0.5cm}
{\Large Trabalho Teorico --- Tema 1\par}
\vspace{0.3cm}
{\large Pontos 1, 2 e 4\par}

\vspace{0.4cm}
{\color{linkazul}\rule{\linewidth}{0.6pt}}

\vfill

\begin{flushright}
{\large\bfseries Autor\par}
{\large Felipe Lazzarini Cunha\par}
\vspace{0.6cm}
{\large\bfseries Grupo 8\par}
\end{flushright}

\vfill
{\large Juiz de Fora\par Julho de 2026\par}
\end{titlepage}

%==============================================================================
% SUMARIO
%==============================================================================
\pagenumbering{arabic}
\tableofcontents
\thispagestyle{fancy}
\newpage

%==============================================================================
% PONTO 1
%==============================================================================
\section{Apresentacao sobre o tema}
\label{sec:apresentacao}

\subsection{Visao geral}

O Tema~1 reune dois sistemas operacionais desenvolvidos pela Google que, embora
compartilhem a mesma mantenedora, representam abordagens arquiteturais
profundamente distintas. De um lado, o \textbf{Android} e o sistema operacional
movel mais difundido do mundo, construido sobre um \emph{kernel} Linux modificado
e amplamente consolidado no mercado~\cite{aosp,androidsite}. De outro, o
\textbf{Fuchsia} e um sistema operacional concebido ``do zero'' pela Google,
estruturado sobre um \emph{microkernel} proprio chamado Zircon e um modelo de
seguranca baseado em capacidades (\emph{capability-based security}), que rompe
com a tradicao dos sistemas derivados do Unix~\cite{fuchsiadev,zircon,
fuchsia_secure}.

Estudar os dois lado a lado e instrutivo justamente pelo contraste: o Android
ilustra a evolucao pragmatica de um \emph{kernel} monolitico maduro adaptado ao
contexto movel, enquanto o Fuchsia exemplifica uma tentativa de repensar, a partir
de premissas modernas de seguranca, modularidade e atualizabilidade, o que deve
ser um sistema operacional de proposito geral~\cite{zircon,fuchsia_components,
fuchsia_updatable}. As secoes seguintes detalham cada sistema; aqui interessa
situar o leitor quanto a origem, ao porte e ao estado atual de ambos.

\subsection{Android}

O Android teve origem na empresa \emph{Android Inc.}, fundada em 2003 e adquirida
pela Google em 2005. Sua primeira versao comercial chegou ao mercado em 2008, com
o lancamento do aparelho HTC Dream~\cite{androidwiki}. Tecnicamente, trata-se de
uma pilha de software cujo nucleo e um \emph{kernel} Linux de suporte de longo
prazo (LTS) acrescido de modificacoes especificas da Google --- os chamados
\emph{Android Common Kernels} (ACKs)~\cite{aospkernel}. Grande parte dessa pilha e
distribuida como software livre por meio do \emph{Android Open Source Project}
(AOSP)~\cite{aosp}.

Em termos de adocao, o Android segue como o sistema operacional de maior alcance
global quando se considera tambem o mercado movel: estatisticas de uso na Web
mantem o Android na lideranca entre sistemas operacionais moveis e tambem no
agregado de todas as classes de dispositivo~\cite{marketshare}. A versao estavel
consolidada mais recente no inicio desta pesquisa era o \textbf{Android~16},
lancado em 10 de junho de 2025; em junho de 2026, a documentacao oficial da
Google ja apresentava o \textbf{Android~17}, ainda com notas de versao beta no
ciclo consultado~\cite{android16,android17}. O alcance do sistema vai de
\emph{smartphones} e \emph{tablets} a televisores, automoveis, dispositivos
vestiveis e experiencias de tela grande baseadas na pilha Android, como o
Googlebook~\cite{googlebook}.

\subsection{Fuchsia}

O Fuchsia veio a publico de maneira atipica: em agosto de 2016, sem qualquer
anuncio formal, jornalistas identificaram no GitHub um repositorio de codigo
mantido pela Google~\cite{fuchsiawiki}. Diferentemente do Android, o Fuchsia nao
se baseia no \emph{kernel} Linux. Seu nucleo --- inicialmente denominado
\emph{Magenta} e renomeado para \textbf{Zircon} em setembro de 2017 --- e um
\emph{microkernel} orientado a objetos com seguranca baseada em capacidades, e o
sistema e implementado em multiplas linguagens, entre elas Rust, C++, C, Dart, Go
e Python~\cite{fuchsiawiki,zircon}.

O Fuchsia chegou a produtos de consumo pela primeira vez em maio de 2021, por meio
de uma atualizacao de software do Google Nest Hub de primeira geracao, que
substituiu o sistema anterior (Cast OS, baseado em Linux); a implantacao se
estendeu aos demais aparelhos de primeira geracao ate agosto de 2021 e alcancou o
Nest Hub de segunda geracao em 2023~\cite{fuchsiawiki}. O projeto, contudo,
enfrentou reducao de equipe a partir de 2023 e, embora siga documentado e
mantido em fontes oficiais recentes, seu papel estrategico tornou-se mais
discreto do que se especulava em seu inicio~\cite{fuchsiawiki,fuchsiadev}.

\subsection{Por que comparar Android e Fuchsia}
\label{sub:porque}

A comparacao entre os dois sistemas e relevante por evidenciar duas filosofias
opostas de construcao de um sistema operacional. O Android adota um \emph{kernel}
monolitico (Linux), no qual grande parte dos servicos --- \emph{drivers}, pilha de
rede, sistemas de arquivos --- executa no espaco de \emph{kernel}; ja o Fuchsia
adota um \emph{microkernel} (Zircon) que mantem no nucleo apenas o essencial,
delegando os demais servicos a componentes isolados em espaco de usuario e
mediados por capacidades~\cite{zircon,fuchsiawiki}. Esse contraste sintetiza um
debate classico de Sistemas Operacionais entre desempenho e simplicidade do modelo
monolitico, de um lado, e isolamento, seguranca e atualizabilidade do modelo
micronucleo, de outro.

Por algum tempo especulou-se que o Fuchsia poderia vir a substituir ou unificar os
sistemas da Google. O movimento recente da empresa, entretanto, aponta em outra
direcao: a documentacao de desenvolvimento para \textbf{Googlebook} apresenta
experiencias de laptop apoiadas na pilha Android e em tecnicas de aplicacoes
adaptativas para telas grandes~\cite{googlebook}. Assim, mais do que disputar o
mesmo nicho, Android e Fuchsia ilustram caminhos tecnologicos distintos que a
Google explorou em paralelo, o que torna o tema especialmente rico para um estudo
de Sistemas Operacionais.

%==============================================================================
% PONTO 2
%==============================================================================
\section{Proposito, evolucao e virtualizacao}
\label{sec:proposito}

\subsection{Proposito de cada sistema}

O \textbf{Android} foi concebido como uma plataforma completa para dispositivos
moveis e embarcados, especialmente aparelhos com tela sensivel ao toque,
conectividade permanente, bateria limitada e forte dependencia de aplicativos de
terceiros. Seu proposito nao e apenas executar processos sobre um \emph{kernel}
Linux, mas oferecer uma pilha integrada que inclui \emph{framework} de aplicacao,
tempo de execucao gerenciado, servicos de midia, seguranca por aplicativo,
distribuicao de software e compatibilidade com grande diversidade de hardware.
Por isso, o sistema combina um nucleo Linux adaptado ao ambiente movel com
mecanismos de isolamento, permissao e atualizacao voltados a um ecossistema amplo
de fabricantes, operadoras, desenvolvedores e usuarios finais~\cite{aosp,
aospkernel,androidsite}.

O \textbf{Fuchsia}, por sua vez, tem um proposito mais experimental e arquitetural:
servir como um sistema operacional de proposito geral, modular, atualizavel e
orientado a seguranca desde a base. Em vez de herdar a arquitetura Unix/Linux, o
Fuchsia e construido sobre o Zircon, um nucleo orientado a objetos que expoe
recursos como capacidades; acima dele, quase todo o software e organizado como
componentes isolados, com dependencias e permissoes declaradas explicitamente
~\cite{zircon,fuchsia_components,fuchsia_secure}. Essa proposta torna o Fuchsia
particularmente interessante para dispositivos conectados e produtos nos quais
isolamento, atualizabilidade e composicao modular sao requisitos centrais, ainda
que sua presenca comercial seja muito mais restrita que a do Android.

\subsection{Evolucao no tempo}

A evolucao do Android pode ser entendida em tres eixos. O primeiro e a expansao
de mercado: desde o Android 1.0, lancado comercialmente em 2008 no HTC Dream, o
sistema deixou de ser uma plataforma de \emph{smartphones} para se tornar uma
familia de variantes usadas tambem em televisores, automoveis, relogios,
\emph{tablets} e outros equipamentos~\cite{androidwiki,androidsite}. O segundo e
a maturacao tecnica da pilha: as versoes iniciais usavam a maquina virtual
Dalvik, enquanto as versoes modernas usam o Android Runtime (ART), que executa o
formato DEX e introduziu compilacao antecipada, melhorias de coleta de lixo e
melhor suporte a depuracao~\cite{androidruntime}. O terceiro eixo e a
modularizacao do proprio sistema. Iniciativas como Treble separaram partes da
plataforma Android das implementacoes especificas de fornecedores, enquanto o
Mainline, introduzido no Android 10, passou a transformar componentes do sistema
em modulos atualizaveis fora do ciclo tradicional de versoes completas
~\cite{treble,androidmainline}.

No Fuchsia, a evolucao foi menos orientada por versoes publicas de produto e mais
por marcos arquiteturais. O projeto apareceu publicamente em 2016, inicialmente
associado ao nome Magenta; em 2017, o nucleo passou a ser conhecido como Zircon
~\cite{fuchsiawiki,zircon}. A primeira implantacao em produto de consumo ocorreu
em 2021, quando o sistema chegou ao Google Nest Hub de primeira geracao, marcando
a transicao do projeto de uma base experimental para uso real em dispositivos
Google~\cite{fuchsiawiki}. Desde entao, o desenvolvimento documentado enfatiza
componentes, pacotes hermeticos, interfaces FIDL, atualizacoes independentes e
compatibilidade com software Linux por meio do Starnix~\cite{fuchsia_components,
fuchsia_updatable,starnix}.

\begin{table}[htbp]
\centering
\caption{Marcos cronologicos de Android e Fuchsia.}
\label{tab:cronologia}
\begin{tabular}{>{\bfseries}l p{5.5cm} p{5.5cm}}
\toprule
Ano & Android & Fuchsia \\
\midrule
2008 & Primeira versao comercial, associada ao HTC Dream. & --- \\
2014--2015 & ART aparece como alternativa ao Dalvik e passa a substitui-lo nas versoes modernas. & --- \\
2016 & --- & Codigo do projeto aparece publicamente, ainda ligado ao nome Magenta. \\
2017 & Treble reorganiza a relacao entre plataforma Android e codigo de fornecedores. & O nucleo passa a ser conhecido como Zircon. \\
2019--2020 & Android 10 introduz Mainline e modulos atualizaveis via APK/APEX. & Consolidacao do modelo de componentes e capacidades. \\
2021 & Android 12 amplia recursos de privacidade e isolamento no sistema. & Primeira implantacao em produto de consumo, no Google Nest Hub. \\
2022--2026 & AVF/pKVM se consolida como mecanismo moderno de virtualizacao protegida em dispositivos ARM64. & Starnix amadurece como camada de compatibilidade para executar programas Linux sem usar o kernel Linux. \\
\bottomrule
\end{tabular}
\end{table}

\subsection{Mecanismos de virtualizacao}

No Android, e importante separar tres ideias que as vezes aparecem misturadas sob
o termo ``virtualizacao''. A primeira e o \textbf{tempo de execucao gerenciado}:
Dalvik e ART nao virtualizam hardware, mas fornecem uma maquina de execucao para
\emph{bytecode} DEX, permitindo que aplicativos escritos em Java ou Kotlin sejam
executados de forma portavel sobre diferentes arquiteturas de CPU
~\cite{androidruntime}. A segunda e o \textbf{isolamento por aplicativo}: cada app
executa em um contexto separado, com identificador, permissoes e politicas de
seguranca proprias; isso e uma forma de compartimentalizacao do sistema, embora
nao seja uma maquina virtual completa.

A virtualizacao de hardware propriamente dita aparece no Android moderno por meio
do \textbf{Android Virtualization Framework} (AVF). Segundo a documentacao do
AOSP, o AVF fornece ambientes de execucao seguros e privados para codigo que
exige isolamento mais forte que o \emph{sandbox} tradicional de aplicativos. Em
dispositivos ARM64 compativeis, o AVF usa o pKVM (\emph{protected kernel-based
virtual machine}) como hipervisor e executa \emph{protected virtual machines}
(pVMs), que podem hospedar, por exemplo, o Microdroid, uma distribuicao Android
minima voltada a execucao isolada~\cite{androidavf}. O ponto central desse modelo
e proteger a memoria e os segredos da maquina virtual mesmo diante de comprometimento
do Android hospedeiro.

No Fuchsia, a arquitetura basica favorece isolamento por desenho, e nao por uma
camada de virtualizacao semelhante ao AVF. O Zircon fornece primitivas de
processos, \emph{threads}, memoria virtual e IPC; acima dele, o \emph{Component
Framework} executa componentes em \emph{sandboxes} com capacidades declaradas
explicitamente~\cite{zircon,fuchsia_components,fuchsia_secure}. Assim, a unidade
normal de isolamento do sistema e o componente com capacidades, nao uma VM.

Para compatibilidade com software Linux, a documentacao atual do Fuchsia destaca
o \textbf{Starnix}. Ele permite executar programas Linux nao modificados em um
sistema Fuchsia, implementando a interface de chamadas de sistema Linux em um
programa de espaco de usuario do Fuchsia, em vez de inicializar um kernel Linux
convidado~\cite{starnix,starnixkernel}. Portanto, para o escopo deste trabalho, a
comparacao mais precisa e: o Android possui virtualizacao protegida de hardware
com AVF/pKVM, enquanto o Fuchsia aposta principalmente em isolamento por
capacidades e componentes, usando Starnix como camada de compatibilidade Linux, e
nao como hipervisor equivalente.

%==============================================================================
% PONTO 4
%==============================================================================
\section{Requisitos do tipo de aplicacao atendido}
\label{sec:requisitos}

\begin{nota}
\textbf{Roteiro (Ponto 4):} Quais os requisitos do tipo de aplicacao que este
SO/kernel atende (baixo tempo de resposta? baixo tempo de retorno? outro?).
\emph{Obs.:} a classificacao do \emph{tipo} de aplicacao (Ponto 3) e o
\emph{como} esses requisitos sao atendidos (Ponto 5) sao de outros integrantes;
aqui foque em \textbf{quais} sao os requisitos.
\end{nota}

\subsection{Requisitos do Android}
% TODO: Baixo tempo de resposta (interatividade, UI a 60/120 Hz),
%       eficiencia energetica (bateria), uso restrito de memoria,
%       responsividade sob restricao de recursos, seguranca/isolamento por app.

\subsection{Requisitos do Fuchsia}
% TODO: Atualizabilidade, seguranca por capacidades, modularidade,
%       escalabilidade entre classes de dispositivo, latencia previsivel.

\subsection{Sintese comparativa dos requisitos}
% TODO: Tabela contrastando requisitos prioritarios de cada sistema.

%==============================================================================
% REFERENCIAS
%==============================================================================
\newpage
\bibliographystyle{plain}
\bibliography{referencias}
% Enquanto a .bib nao estiver populada, pode-se comentar as duas linhas acima
% e usar a lista manual abaixo:
% \begin{thebibliography}{9}
% \bibitem{android} Android Open Source Project. \url{https://source.android.com/}.
% \bibitem{fuchsia} Fuchsia. \url{https://fuchsia.dev/}.
% \end{thebibliography}

\end{document}
